%! Author = lili
%! Date = 7/24/26

\subsection{Zusammenhang Poissonscher Grenzwertsatz, Zentraler Grenzwertsatz, Satz von de Moivre-Laplace und Gesetz der großen Zahlen}\label{subsec:zusammenhang-poissonscher-grenzwertsatz-zentraler-grenzwertsatz-satz-von-de-moivre-laplace-und-gesetz-der-groen-zahlen}
\index{Gesetz der großen Zahlen} \index{Poissonscher Grenzwertsatz} \index{Moivre-Laplace}
\index{Zentraler Grenzwertsatz}

\begin{mytable}
    \begin{tabular}{|l|l|l|}
        \hline
        \textbf{Satz} & \textbf{Fragestellung} & \textbf{Ergebnis} \\
        \hline
        Gesetz der großen Zahlen & Konvergiert der Durchschnitt? &
        Der Stichprobenmittelwert nähert sich dem Erwartungswert. \\
        \hline
        Zentraler Grenzwertsatz &
        Wie schwanken Summen bzw. Mittelwerte? &
        Die normierte Summe wird asymptotisch normalverteilt. \\
        \hline
        Satz von de Moivre-Laplace &
        Spezialfall des ZGWS &
        Die Binomialverteilung wird durch eine Normalverteilung angenähert. \\
        \hline
        Poissonscher Grenzwertsatz &
        Seltene Ereignisse &
        Die Binomialverteilung wird durch eine Poissonverteilung angenähert. \\
        \hline
    \end{tabular}
\end{mytable}


\subsubsection*{Gesetz der großen Zahlen}

Seien $X_1,\dots,X_n$ unabhängige und identisch verteilte Zufallsvariablen mit Erwartungswert $\wfk$. Dann gilt

\[
    \overline{X}_n=\frac1n\sum_{i=1}^n X_i.
\]

Nach dem (schwachen) Gesetz der großen Zahlen gilt

\[
    \overline{X}_n \xrightarrow{P} \wfk.
\]

Beim starken Gesetz der großen Zahlen gilt sogar

\[
    \overline{X}_n \xrightarrow{\text{f.s.}} \wfk.
\]

\paragraph{Interpretation:}

Mit wachsender Stichprobengröße nähert sich der Stichprobenmittelwert immer stärker dem Erwartungswert an.

Das Gesetz der großen Zahlen macht jedoch keine Aussage darüber, wie groß die zufälligen Abweichungen sind oder wie diese verteilt sind.

\subsubsection*{Zentraler Grenzwertsatz}

Unter denselben Voraussetzungen sowie einer endlichen Varianz $\sigma^2$ gilt

\[
    \frac{\sum_{i=1}^n X_i-n\mu}{\sigma\sqrt n}
    \Longrightarrow N(0,1).
\]

Äquivalent für den Stichprobenmittelwert:

\[
    \frac{\overline X_n-\mu}{\sigma/\sqrt n}
    \Longrightarrow N(0,1).
\]

\paragraph{Interpretation}

Die zufälligen Abweichungen des Mittelwertes vom Erwartungswert sind für große Stichproben näherungsweise normalverteilt.

Der Zentrale Grenzwertsatz liefert somit deutlich mehr Information als das Gesetz der großen Zahlen.


\subsubsection*{Zusammenhang zwischen GGZ und ZGWS}

Der Zentrale Grenzwertsatz zeigt, dass

\[
    \overline X_n-\mu
    =
    O\!\left(\frac1{\sqrt n}\right).
\]

Da

\[
    \frac1{\sqrt n}\longrightarrow 0,
\]

folgt daraus unmittelbar die Aussage des schwachen Gesetzes der großen Zahlen.

Anschaulich gilt:

\begin{itemize}
    \item Das Gesetz der großen Zahlen beschreibt, dass der Fehler gegen Null verschwindet.
    \item Der Zentrale Grenzwertsatz beschreibt zusätzlich die Verteilung dieses Fehlers.
\end{itemize}


\subsubsection*{Satz von de Moivre-Laplace}

Sei

\[
    X_n\sim\operatorname{Bin}(n,p).
\]

Dann gilt

\[
    \frac{X_n-np}{\sqrt{np(1-p)}}
    \Longrightarrow N(0,1).
\]

Der Satz von de Moivre-Laplace ist somit ein Spezialfall des Zentralen Grenzwertsatzes für Bernoulli-Zufallsvariablen.

Er liefert die Approximation

\[
    \operatorname{Bin}(n,p)
    \approx
    N(np,np(1-p)).
\]

Diese Approximation ist insbesondere geeignet, wenn

\begin{itemize}
    \item $n$ groß ist,
    \item $p$ weder sehr klein noch sehr groß ist.
\end{itemize}

\subsubsection*{Poissonscher Grenzwertsatz}

Betrachte nun

\[
    X_n\sim\operatorname{Bin}(n,p_n),
\]

wobei

\[
    p_n\longrightarrow 0,
    \qquad
    np_n\longrightarrow\lambda.
\]

Dann gilt

\[
    X_n
    \Longrightarrow
    \operatorname{Pois}(\lambda).
\]

Hier wird die Binomialverteilung nicht durch eine Normalverteilung, sondern durch eine Poissonverteilung angenähert.

Diese Approximation eignet sich insbesondere für seltene Ereignisse, beispielsweise:

\begin{itemize}
    \item Defekte Bauteile,
    \item Verkehrsunfälle,
    \item Druckfehler.
\end{itemize}

\subsubsection*{Verhältnis zwischen de Moivre-Laplace und Poissonschem Grenzwertsatz}

Beide Grenzwertsätze gehen von der Binomialverteilung aus, betrachten jedoch unterschiedliche Grenzfälle.

\[
    \boxed{
        \begin{array}{ccc}
            \operatorname{Bin}(n,p)
            & &
            \operatorname{Bin}(n,p_n)
            \\[0.3cm]
            p=\text{konstant}
            &&
            p_n\to0,\;np_n\to\lambda
            \\[0.3cm]
            \Downarrow
            &&
            \Downarrow
            \\[0.3cm]
            N(np,np(1-p))
            &&
            \operatorname{Pois}(\lambda)
        \end{array}
    }
\]

\subsubsection*{Gesamtübersicht}

\[
    \boxed{
        \begin{array}{c}
            \text{Beliebige i.i.d. Zufallsvariablen}
            \\[0.4cm]
            \Downarrow
            \\[0.4cm]
            \text{Gesetz der großen Zahlen}
            \\
            \overline X_n\to\mu
            \\[0.5cm]
            \Downarrow
            \\[0.5cm]
            \text{Zentraler Grenzwertsatz}
            \\
            \dfrac{\overline X_n-\mu}{\sigma/\sqrt n}
            \Longrightarrow N(0,1)
            \\[0.5cm]
            \Downarrow
            \\[0.5cm]
            \text{de Moivre-Laplace}
            \\
            \text{(Spezialfall für Bernoulli- bzw. Binomialverteilungen)}
        \end{array}
    }
\]

\paragraph{Merksätze}
\begin{itemize}
    \item \textbf{Gesetz der großen Zahlen:} Der Stichprobenmittelwert konvergiert gegen den Erwartungswert.
    \item \textbf{Zentraler Grenzwertsatz:} Die normierten Abweichungen vom Erwartungswert sind asymptotisch normalverteilt.
    \item \textbf{de Moivre-Laplace:} Spezialfall des Zentralen Grenzwertsatzes für Binomialverteilungen.
    \item \textbf{Poissonscher Grenzwertsatz:} Spezialfall für seltene Ereignisse mit $p\to0$ und $np\to\lambda$; die Grenzverteilung ist eine Poissonverteilung.
\end{itemize}
